\section{Введение}

Число параметров описывает размер представления нейросети, но само по себе не
определяет стоимость передачи её предсказаний или качество обобщения. Один и тот
же предиктор можно представить разными наборами весов; одна архитектура при
разных данных и правилах обучения даёт разные результаты. Принцип минимальной
длины описания (minimum description length, MDL) предлагает выбирать модель
вместе с кодом, который объясняет наблюдения коротким сообщением
\citep{grunwald2004}. Для нейросетей этот подход реализуется через вариационное
кодирование, приближения маргинального правдоподобия, последовательные коды и
компрессию весов \citep{graves2011,ritter2018,blier2018,arora2018}.

Однако даже длина в битах не является универсальной характеристикой модели:
результат зависит от того, \emph{что} восстанавливает получатель. В
prequential-кодировании передаются наблюдения, и декодер обучается на уже
раскрытом префиксе \citep{blier2018,bornschein2022}. В requential-кодировании
передаются индексы примеров, сгенерированных из распределения ученика и
выбранных учителем \citep{finzi2026,qiu2026}. В коде весов передаётся конечное
состояние модели. Обучающие последовательности с совершенно разными длинами
могут приводить к близким конечным состояниям. Поэтому сравнение этих величин
требует явного задания протокола.

Мы исследуем асимптотику стоимости requential-кодирования при заданной
точности конечного предиктора. Рассматриваются три вопроса. Как выбирать моменты
обновления ученика, если каждый переданный пакет имеет собственные накладные
расходы? Определяется ли накопленный риск геометрическим коэффициентом
байесовского обучения, когда вместо интегрирования по параметрам используется
градиентный спуск? И когда полную запись обучения можно заменить коротким
описанием его конечного результата?

Мы получаем следующие результаты.
\begin{enumerate}
\item Для замороженного внутри пакета предиктора с априорным распределением
$\Beta(a,b)$ и нулевого
учителя доказываем
\[
 C_{\rm fr}(\varepsilon)=A(a)\ln(1/\varepsilon)+O(1),\qquad
 A(a)=\min_{r>1}\frac{L(a(r-1))}{\ln r},
\]
где $L(h)$ --- ожидаемая длина $\delta$-кода Элиаса для геометрического индекса
со средним $e^h$. Параметр $b$ влияет на конечные поправки, но не на ведущий
коэффициент. При ограничении среднего числа кандидатов область минимизации
по $r$ сужается явным образом.
\item Показываем, что совместное байесовское предсказание пакета меняет ведущий
коэффициент на $a/\ln 2<A(a)$. Таким образом, при фиксированных априорном
распределении, учителе и конечном риске стоимость зависит от способа
формирования пакета.
\item Для точного градиентного спуска в факторизованной регрессии с нулевой
истинной матрицей доказываем степенной закон, зависящий от предельного
дисбаланса норм. Дополняем его случаем точного баланса:
$K_t\sim(8\gamma^2\ln^2t)^{-1}$ при шаге $\gamma/(t+t_0)$.
\item Для сжимающих обновлений строим requential-описание конечной точки по
короткому суффиксу меток. Его ожидаемая длина ограничена через точность
восстановления и не растёт с длиной исходной траектории.
\end{enumerate}

Эти результаты относятся к явно заданным классам кодов и алгоритмов. В
частности, нижняя оценка для протокола с обязательным индексом на пакет не
является нижней оценкой колмогоровской сложности конечной модели. Аналитические
примеры ниже служат проверкой формул; перенос выводов на большие нейросети
требует отдельных экспериментов.
